\section*{Part VII: Financing Physical Expansion}
\addcontentsline{toc}{section}{Part VII: Financing Physical Expansion}


\subsection{54. Equity, debt, grants and property partners}

Potential sources include retained earnings, owner/sponsor capital, restricted construction philanthropy, equipment grants, bank mortgage or term debt, lease financing, vendor credit, landowner joint development, university/industry co-investment and project vehicles. Each has different claims on property, cash flow, control and exit. VC investment into an education software company is not automatically appropriate financing for a charitable academic campus.

Do not pledge an operating company's most valuable assets to support a campus whose business model cannot service the debt without a well-understood risk decision.


\subsection{55. Why project finance is not simply a loan secured by property}

True project finance usually depends heavily on project cash flows, security, risk allocation and a ring-fenced structure; ordinary corporate loans may use broad sponsor recourse, guarantees or cross-collateral. A university teaching building without contracted stable cash flows may be difficult to finance as a standalone cash-generating project.

Ask whether cash comes from rents, tuition, contracted facility use, institutional payments or donor sources. \textbf{Repayment needs predictable cash flows, not merely appraised asset value.}


\subsection{56. Debt-service stress test}

\textbf{[Model]} A \textnormal{₹}10-crore loan amortized by equal annual payments over 10 years at 11% hypothetical annual interest requires approximately \textbf{\textnormal{₹}1.698 crore per year}. If a lender or internal policy seeks a DSCR of 1.30 under its defined CFADS, the project needs roughly \textbf{\textnormal{₹}2.207 crore annual cash flow available for debt service} at base case, with more cushion for downside. Actual lender amortization, grace period, fees and covenants will differ.

Now reduce realized tuition/campus revenue by 25%, increase operating cost by 15%, delay occupancy 12 months and raise rates. If repayment becomes impossible under plausible shocks, reduce debt, build in phases or wait.


\subsection{57. Debt security and hidden founder exposure}

Review collateral, pledge of shares, personal guarantees, cross-default clauses, cash-sweep provisions, debt-service reserves, financial covenants, mortgage registration, restricted payments and lender step-in rights. A facility loan may seem non-dilutive while putting founder personal assets and the operating company at risk.

An independent lawyer and finance adviser should prepare a \textbf{recourse map} showing exactly which person and asset is exposed under default. Ownership preservation is not valuable if 1 missed repayment can destroy the entire venture.


\subsection{58. Lender drawdown versus capital call sequencing}

Construction lenders may require sponsor equity to be contributed 1st, or impose draw conditions tied to permits, valuations, architect certificates and progress. Philanthropic commitments may have their own matching conditions. Plan the order, not only the total. A failure to meet 1 condition can freeze all downstream funding while contractors remain due.

Require at least 1 dedicated project finance controller once commitments become material. The routine enterprise accounts team may not be sufficient for a complex build.


\subsection{59. The restricted-funds problem}

Donor capital may be restricted to land, a particular building, equipment, scholarships, or an endowment. It may not lawfully or contractually support construction interest, overruns or payroll. A binding pledge can still be payable over years. Maintain a restricted-fund ledger and permission-to-spend review by an independent finance officer.

Do not transfer donated educational property to a founder-controlled operating company without the legal authority, fair valuation, conflict controls and donor consent that the circumstances require.


\subsection{60. Capital recycling and leased assets}

Instead of acquiring every adjacent function, the institute may lease housing, partner for sports facilities, use nearby private student accommodation, rent overflow labs, or outsource food/maintenance services under quality and safeguarding contracts. This can lower initial capital but requires careful vendor oversight and contingency if a provider fails.

The institution's strategic identity should derive from teaching and research standards, not from owning every service provider in its supply chain.
